\documentclass[../../../thesis.tex]{subfiles}
\begin{document}
\subsection{Hlavná textová časť}
Samotný autorský obsah práce začína až tu.
Tradične text členíme na úvod, jadro a~záver,
pričom úvod a~záver sú samostatné kapitoly,
ktoré nečíslujeme a~je vhodné,
ak ich označíme nadpismi \emph{Úvod} a~\emph{Záver}.
Strednú časť -- jadro -- neozначаujeme.

\subsubsection{Úvod}
Prvá kapitola hlavnej časti práce má názov úvod, nečíslujeme ju.
Ide o~ucelený text v~rozsahu niekoľkých súvislých odsekov textu,
v~ktorých stručne a~výstižne charakterizujeme stav poznania
a~praxe v~danej oblasti,
oboznámime čitateľa s~cieľmi a~závermi práce.
Nosnou myšlienkou úvodu okrem uvedenia čitateľa do problematiky
je jasná motivácia autora a~jeho postoje,
ktoré viedli k~spracovaniu témy práce [3].

Nepísané pravidlo hovorí,
že úvod a~záver práce sa píšu až ako posledné.
Tento poznatok vyplýva z~praxe a~má dva dôvody:
1.~na začiatku nemusí byť úplne zrejmé, čo všetko sa v~práci naozaj objaví;
2.~úvod predstavuje samostatnú literárnu formu,
na ktorej sa neskúsený autor zasekne už na začiatku.
Aby sme sa tomu vyhli,
necháme si jeho napísanie až na záver,
keď už bude väčšina hlavného obsahu práce hotová.

Text úvodu sa nachádza v súbore \verb|includes/introduction.tex|
a~jeho sadzbu zariadi makro \verb|\FEIintroduction{includes/introduction}|.

\subsubsection{Jadro}
Táto časť práce \emph{nezačína} nadpisom \emph{Jadro}.
Obsah jadra členíme zvyčajne na niekoľko číslovaných
kapitol počínajúc číslom 1.
Prvá kapitola býva prehľad súčasného stavu problematiky,
ale môže mať aj iný názov,
napríklad \emph{Teoretická časť,} alebo rovno názov oblasti,
o~ktorej sa v nej bude písať
(trebárs \emph{Metóda prenosových matíc}).

Pri písaní strednej časti práce nemusíme postupovať úplne
striktne podľa tohto návodu.
Treba však pamätať na to,
aby sme jasne oddelili poznatky,
ktoré pochádzajú od iných autorov,
a~sú súčasťou všeobecného prehľadu,
od poznatkov a~výsledkov samotnej práce autora.
Nemusia byť oddelené fyzicky v~rôznych odsekoch,
či kapitolách, z~textu však musí byť jasné,
ktoré výsledky sú originálne a~ktoré sú prebrané.
Odporúčaná štruktúra tejto časti je na
strane~\pageref{sec:StrukturaPrace}.

Samotný obsah jadra sa nachádza v~súbore \verb|includes/core.tex|.
Do hlavného dokumentu \verb|thesis.tex| ho načíta makro \verb|\FEIcore{includes/core}|.
Parameter makra je názov súboru bez prípony.
Ak je \verb|core.tex| príliš obsiahly,
môžeme jednotlivé kapitoly uložiť do samostatných
súborov a tie načítať do \verb|core.tex|
pomocou \TeX-ového príkazu \verb|\input|.

\subsubsection*{\normalsize Súčasný stav riešenej problematiky
doma a~v~zahraničí}
Podľa zvyklostí by malo približne 30\,\% práce obsahovať prehľad
súčasného stavu a~poznatkov v~oblasti,
ktorej sa týka predkladaná práca.
Ide o~veľmi dôležitý aspekt,
ktorým študent preukáže,
že je schopný problematiku naštudovať,
porozumieť jej a~napísať o~nej súvisly text.
Dokáže na základe existujúcich poznatkov vysvetliť javy,
ktoré v~práci študuje.

Kľúčová činnosť pri príprave textu je štúdium prác publikovaných
u~nás a~v~zahraničí.
Nejde iba o~to, že autor píše myšlienky, ktoré sa kdesi dozvedel,
mal by tiež poznať ich primárne zdroje,
správne s~nimi pracovať a~citovať ich.
Dôležitý prínos študenta spočíva v~spájaní viacerých poznatkov
z~rôznych zdrojov do nového celku.

\subsubsection*{\normalsize Cieľ práce}
Bakalárska a diplomová práca má jasne uvedené ciele v zadaní práce. Nie je preto nutné uvádzať samostatnú kapitolu, kde budú ciele ešte raz vymenované. Je však žiadúce, ak sa zmienka o jednotlivých cieľoch v texte vyskytuje a poukazuje sa na ich splnenie, nesplnenie, prípadne ak hlavné ciele pozostávajú z čiastkových cieľov, treba ich jasne špecifikovať.

\subsubsection*{\normalsize Metodika práce a metódy skúmania}
V experimentálnych prácach býva v tejto časti podrobne zdokumentované prístrojové vybavenie, riadiaci a simulačný softvér, laboratórne podmienky a podobne. Metodické usmernenie [3] odporúča nasledujúci obsah tejto časti práce: a) charakteristika objektu skúmania, b) pracovné postupy, c) spôsob získavania údajov a ich zdroje, d) použité metódy vyhodnotenia a interpretácie výsledkov, e) štatistické metódy.

\subsubsection*{\normalsize Výsledky práce a diskusia}
Študent zaujme k získaným výsledkom jasné postoje,
porovnáva ich s inými autormi, prípadne navrhuje ich ďalšie aplikácie.
Zhodnotí a~komentuje ich na základe štatistického spracovania dát (smerodajné odchýlky, priemery, regresie a podobne).
Odporúčame, aby táto časť tvorila 30 až 40 percent záverečnej práce.
Môžeme ju rozdeliť na dve samostatné podkapitoly: sumarizáciu výsledkov a~diskusiu formou eseje.

\subsubsection{Záver}
Záver práce predstavuje samostatnú nečíslovanú kapitolu
v~rozsahu niekoľkých odsekov alebo strán.
Obsahuje zhrnutie výsledkov vo vzťahu k~stanoveným cieľom~[3].
Rovnako, ako pri úvode, treba si dať
aj na kompozícii záveru zvlášť záležať.
Väčšina čitateľov si prečíta v~prvom rade úvod a~záver práce,
aby zistili, či im stojí za to pustiť sa do podrobnejšieho
štúdia celého textu.
Aj oponent vychádza najmä z dobre spracovaného záveru.

Jasne deklarujeme splnenia cieľov a naznačíme ďalšie možné smerovanie študovanej problematiky. Vyjadrujeme sa pozitívne. Ak sa nepodarilo úplne naplniť niektorú z~pôvodných predstáv, nerozpisujeme sa o tom.

Ako príklad použijeme nepríjemnú modelovú situáciu,
ktorá môže počas výskumu nastať.
Povedzme, že cieľ záverečnej práce bol odmerať
optické parametre tenkých TiO$_2$ vrstiev.%
\footnote{TiO$_2$ je chemická značka oxidu titaničitého,
ktorý sa používa napríklad pri solárnych článkoch
ako priehľadná vrchná elektróda.
Ide totiž o~typ oxidu s~vlastnosťami polovodičov,
čiže môže za určitých podmienok viesť elektrický prúd.
Zároveň je pre viditeľné svetlo priehľadný,
čo nebýva pri polovodičoch bežné.
Optické a~elektrické vlastnosti vrstvy TiO$_2$
často závisia od parametrov technologického procesu.}
Z~dôvodu havárie zariadenia sa nepodarilo takéto vzorky získať
a~v~skutočnosti sme mohli pracovať iba
s~tradičnými SiO$_2$ vrstvami.%
\footnote{Oxid kremičitý sa v~mikroelektronike používa
ako nevodivá izolačná vrstva.
Jeho materiálové vlastnosti sú veľmi dobre preskúmané
a~všeobecne známe.
S~jeho amorfnou formou sa v~každodennom živote bežne stretávame,
je to obyčajné sklo.}
Vzniknutú situáciu zhodnotíme v~závere vecne a~pravdivo:

\begin{quote}
  \emph{Aj napriek poruche technologického zariadenia sme
  dokázali zabezpečiť náhradné vzorky a realizovať merania
  optických vlastností tenkých vrstiev termálneho
  SiO$_\mathit{2}$.
  Poznatky, ktoré sme získali pri práci s~pokročilými
  experimentálnymi zariadeniami následne využijeme vo výskume
  materiálových vlastností TiO$_\mathit{2}$ vrstiev.
  V~diskusii sme naznačili možné rozšírenie existujúcich
  metód na tento druh materiálu.}
\end{quote}

Ak priznáme, že zariadenie sa pokazilo
a tým pádom sme nesplnili ciele,
stane sa záverečná práca neobhájiteľnou.
Nasledujúci príklad je ukážka takejto nevhodnej formulácie:

\begin{quote}
  \emph{Počas prípravy tenkých vrstiev došlo k neočakávanej
  poruche technologického zariadenia,
  ktorá znemožnila výrobu plánovaných vzoriek.
  Merania optických parametrov TIO$_\mathit2$
  sme preto nerealizovali.
  Veríme, že experimenty s~náhradnými vzorkami tenkých vrstiev
  termálneho SiO$_\mathit2$ pomôžu v~budúcnosti
  aj pri výskume iných materiálov.}
\end{quote}

Text obsahuje tri zápory, je pesimistický,
s~nejasným výhľadom do budúcnosti.
Cítiť z~neho sklamanie a~frustráciu zo vzniknutej situácie,
ktorá sa javí ako neriešiteľná.
Jednoznačne sme priznali nesplnenie cieľa.
Aj keď sme urobili úspešné náhradné merania,
z~textu to nie je zrejmé.
Záverečné tvrdenie o~možnosti využitia výsledkov
v~sebe navonok ukrýva istú nádej,
v~skutočnosti však iba potvrdzuje to,
že chceme mať toto fiasko čím skôr za sebou.

Pozor ale aj na prílišnú pozitivitu.
Tá môže, paradoxne, nedostatky ešte viac zvýrazniť.
Nasledujúca ukážka je síce optimistická,
avšak do textu práce taktiež nevhodná:

\begin{quote}
  \emph{Vďaka drobnej poruche technologického zariadenia sme
  mohli realizovať merania optických vlastností tenkých vrstiev
  termálneho SiO$_\mathit2$ a~získať tak unikátne výsledky.
  Nesmierne bohaté skúsenosti s~najkvalitnejšími meracími
  aparatúrami využijeme aj v~nadväzujúcom výskume.
  Rozšírenie nadobudnutých kompetencií na iné materiály
  považujeme za najväčší prínos predkladanej práce.}
\end{quote}

V tomto príklade vidieť prílišnú snahu zahladiť škody
a~vychvaľovať sa výsledkami,
ktoré v~skutočnosti nemajú zvláštny význam.
Je totiž málo pravdepodobné,
aby s~SiO$_2$ vznikli unikátne výsledky.
Text obsahuje nevhodné absolútne kvantifikátory
(\emph{nesmierne bohaté skúsenosti, najkvalitnejšie aparatúry,
najväčší prínos});
bagatelizuje nehodu, dokonca jej ďakuje
(\emph{vďaka drobnej poruche}),
čím na ňu zbytočne upozorňuje;
zámerne sa nezmieňuje o~pôvodných TiO$_2$ vrstvách.
Nadužívaním cudzích slov (\emph{kompetencie})
autori zväčša maskujú rôzne nedostatky,
napríklad vlastnú neistotu.

Zapamätáme si, že vedecký text musí byť jasný, pravdivý a vecný.
Očistíme ho od akýchkoľvek citových výlevov v prvom rade tým,
že sa vyhýbame extrémnym kvantifikátorom. Nepoužívame ani tieto:
\emph{všetci, nikdy, žiaden, každý jeden,}
pokiaľ nepíšeme matematické vety alebo logické výrazy.
Ak sa napríklad nepodarilo naprogramovať ani jeden fungujúci kód,
nenapíšeme,
že \emph{žiaden program, ktorý sme sa snažili vytvoriť nefunguje}.
Povieme to miernejšie: \emph{snaha o~vytvorenie funkčného
programu viedla k~menej presvedčivým výsledkom}.
Negatívnu skutočnosť formulujeme pozitívne.

Ani pozitívne prínosy zbytočne nepreceňujeme.
Necháme ich, nech sa chvália samé.
Namiesto prehnaného zdôrazňovania:
\emph{Úžasné výsledky všetkých meraní sme dosiahli
vďaka perfektne pripraveným vzorkám},
napíšeme vecne:
\emph{Jednotlivé merania boli úspešné aj
vďaka kvalitným vzorkám.}

Súbor so záverom v~priečinku \verb|includes| má
názov \verb|conclusion.tex|
a~do dokumentu sa dostane prostredníctvom makra
\verb|\FEIconclusion{includes/conclusion}|
v~hlavnom súbore projektu \verb|thesis.tex|.

\subsubsection{Zoznam použitej literatúry}
Citované zdroje označujeme v texte číslom v hranatých zátvorkách. Ide o poradové číslo uvedenia publikácií tak, ako sa postupne s nimi v texte pracuje.

Po kapitole \emph{Záver} nasleduje ďalšia nečíslovaná kapitola
s~názvom \emph{Literatúra},
ktorá obsahuje číslovaný zoznam všetkých
citovaných literárnych zdrojov v spomínanom poradí.
Forma tohto zoznamu je pomerne komplikovaná a~podrobne
ju opisuje norma ISO 690: 2023 Dokumentácia -- Bibliografické odkazy -- Obsah, forma a štruktúra \cite{iso690}.
Citovanie je v~\LaTeX-u vynikajúco vyriešené.
V~tomto dokumente citujeme pomocou nadstavby Bib\LaTeX.
Podrobne sa citáciám budeme venovať v~\ref{sec:citation}. kapitole.
\end{document}